\section{相关工作}
\label{sec:related}

本节按三条技术线索回顾与本文最直接相关的研究，
并在每条线索末尾点明 DASD-PPO 与现有方法的关键差异。

\subsection{TDMA 调度与空间复用}

固定 TDMA 通过预分配所有者时隙保证无冲突接入，
但在动态拓扑下大量所有者时隙闲置~\cite{author2024survey}。
为提升频谱效率，
图染色与几何冲突图类方法（如 STDMA~\cite{stdma}）
将时隙复用建模为染色问题，
通过两跳安全约束允许相距足够远的节点同时占用同一时隙。
\textbf{Z-MAC}~\cite{zmac} 采用 CSMA/TDMA 混合机制：
非所有者节点在所有者未占用时通过竞争争抢时隙，
实现负载自适应；
\textbf{TRAMA}~\cite{trama} 通过基于优先级的确定性两跳仲裁完成无冲突调度，
并为节能引入信令调度树。
\emph{空间复用与吞吐量关系}的理论分析进一步给出了密集拓扑下可达吞吐的上界~\cite{spatial_reuse}。
然而，上述方法的调度规则结构与关键参数均在部署前离线确定、
运行期不随业务到达分布与拓扑变化做在线自适应，
因此在拓扑突变与到达过程剧烈波动场景下仍难以兼顾吞吐与长尾时延。
本文不试图替换其中任意一种协议状态机，
而是在 PPO 学习框架内显式引入\emph{两跳安全掩膜}与\emph{服务债务}两项归纳偏置，
将传统调度的拓扑感知先验作为可学习策略的边界约束。

\subsection{强化学习驱动的 MAC 调度}

用学习方法求解 MAC 决策已积累了若干代表性工作。
其中 DLMA~\cite{drlmac2019} 与 PPOMA~\cite{ppoma2023} 面向单节点共享一条异构信道的二元接入决策，
分别以 DQN 与 PPO 为策略骨干；
Dutta 等~\cite{bandit} 则将每节点的时隙选择建模为分布式多臂老虎机，
作为非深度学习的轻量替代
（MAB 在更广义无线决策中的综述可见~\cite{mab_5g}）。
本文沿用其中 PPO~\cite{schulman2017ppo}、GAE~\cite{schulman2016gae} 与 LSTM~\cite{lSTM_seq} 等成熟组件，
但把目标场景从"单节点二元接入"扩展为"多节点 \(\times\,M\) 维时隙申请 + 空间复用 + 公平性"
的分布式 TDMA 调度。

然而，DLMA/PPOMA 类方法面向的均是\emph{单节点 \(\times\) 单信道}的接入博弈，
其策略空间仅含"发/不发"二元动作，
与本文关注的\emph{分布式多节点 \(\times M\) 维时隙申请 + 空间复用 + 公平性}调度问题在尺度与结构上均不重合。
据我们所知，
将 PPO 策略与完整三阶段 TDMA 调度框架、空间复用约束与服务债务机制\emph{统一在同一在线学习管线}下的代表性公开实现仍属空白。
直接把上述二元决策方法移植到多时隙申请场景时，
还会暴露两类共性缺陷：
\emph{其一}，由策略直接输出申请概率会覆盖已有启发式调度器的退避机制与安全先验，
易在训练初期破坏可行域、引发过申请；
\emph{其二}，\(M\) 维 Bernoulli 申请动作只对应一个聚合奖励，
单时隙决策无法获得精确信用归因，
PPO 的策略熵在小批量更新下容易塌缩。
DASD-PPO 通过\emph{乘数模式策略合成}（保留启发式作为先验基准）与\emph{逐时隙信用分配}
（Critic 输出 \(M\) 维价值、PPO 损失逐时隙独立计算）同时回应上述两个缺陷。

\subsection{公平性与队列稳定性}

公平性是多用户共享信道场景中与吞吐量同等重要的指标。
Jain 等~\cite{jain1984} 定义的 Jain 公平指数已成为评价资源分配公平性的标准工具，
本文沿用其在最后 100 帧的局部值与全过程趋势两种统计。
\textbf{Lyapunov drift 方法}~\cite{neely2010stochastic} 提供了在保证队列稳定的同时最大化效用的形式化框架，
但其拉格朗日乘子需要在线估计且对环境噪声敏感。
本文在工程实验中也曾尝试以拉格朗日变体显式约束饥饿，
然其约束信号在仿真中容易饱和到 \(1.0\)（见 \Cref{sec:discussion}）；
DASD-PPO 选择以\emph{服务债务}的二值指示量驱动确定性动作推力，
辅以请求预算上限抑制过申请，
在工程实现上更稳定，
同时保留了 Lyapunov 思想中"将公平性作为可调节项叠加到效用最大化"的结构。
